% CREATED BY MAGNUS GUSTAVER, 2020
\chapter{Methods}
\label{chap:methods}

This chapter describes NanoMem as a query-time memory evidence provider. The
chapter starts by formalizing the task as evidence-state construction rather
than session retrieval. It then describes the system architecture, the Temporal
Evidence Pool, temporal normalization, the Planner--Retriever--Synthesizer
loop, the training objective, the reward design, the ChainMem benchmark, and the
implementation boundary. The common thread is the separation between memory and
answering: NanoMem constructs compact, source-grounded evidence and a
sufficiency verdict, while a downstream answer model consumes that evidence to
produce the final answer.

\section{Problem Formulation}
\label{sec:methods-formulation}

Let the long-term conversational memory bank be
\[
  \mathcal{M} = [(t_1,S_1),(t_2,S_2),\ldots,(t_N,S_N)],
\]
where $S_i$ is a historical session and $t_i$ is the timestamp at which that
session was recorded. Each session contains an ordered sequence of conversation
units,
\[
  S_i = \{c_{i1},c_{i2},\ldots,c_{iL_i}\},
\]
where a unit may be a message, turn, or chunk. Given a user query $q$, a
standard memory retrieval system tries to select a subset of sessions
$\mathcal{S}\subseteq\mathcal{M}$ and provide them to an answer model. This
formulation is insufficient for the evidence addressability gap introduced in
Chapter~\ref{chap:introduction}: the query may not contain the cue needed to
retrieve the answer-bearing session.

NanoMem instead treats the task as construction of an evidence state. The
system incrementally builds a Temporal Evidence Pool $H_T$ such that a
downstream answer model can answer from $(q,H_T)$ without reading the complete
history. Each evidence item is query conditioned. A session can support
different events depending on the query, so the system does not assume that all
useful event interpretations can be fixed at ingestion time.

The formulation distinguishes three time variables. The \emph{session time}
$t_i$ is the timestamp of the source session. The \emph{event time} is the time
or interval when the event described by an evidence record actually happened or
was valid. The \emph{query time} is the timestamp that anchors relative temporal
expressions in the current query. These variables often differ. If a user says
on April 2 that they travelled to Europe last week, the session time is April 2,
while the event time is the preceding week.

\begin{figure}[H]
  \centering
  \fbox{\begin{minipage}{0.92\textwidth}
    \centering
    \vspace{0.8em}
    \textcolor{red}{TODO: Replace this placeholder with the planned task
    redefinition diagram. The left side should show conventional session
    retrieval returning raw sessions. The right side should show NanoMem
    constructing source-grounded events with session time, event time, and a
    verdict.}
    \vspace{0.8em}
  \end{minipage}}
  \caption{From session retrieval to evidence-state construction. NanoMem
  changes the memory objective from selecting sessions to constructing a
  structured Temporal Evidence Pool for a downstream answer model.}
  \label{fig:methods-task}
\end{figure}

\begin{table}[H]
  \centering
  \caption{Core notation and fields used in the NanoMem method.}
  \label{tab:methods-symbols}
  \scriptsize
  \begin{tabular}{p{0.18\textwidth}p{0.34\textwidth}p{0.18\textwidth}p{0.20\textwidth}}
    \hline\hline
    Symbol or field & Meaning & Source & Role in the system \\
    \hline
    $\mathcal{M}$ & Multi-session memory bank & Stored dialogue history & Search substrate \\
    $S_i$ & Historical session containing turns or chunks & Memory database & Retrieval unit \\
    $t_i$ & Timestamp of session $S_i$ & Session metadata & Session-time grounding \\
    $q$ & Current user query & User input & Defines query-conditioned evidence need \\
    $H_t$ & Temporal Evidence Pool after round $t$ & Synthesizer output & Working evidence state \\
    \texttt{session\_id} & Identifier of the source session & Retriever/session metadata & Source attribution \\
    \texttt{session\_time} & Timestamp of the cited session & Session metadata & Mention-time record \\
    \texttt{event\_time} & Time or interval of the described event & Temporal normalization or Synthesizer reasoning & Retrieval cue and evidence field \\
    \texttt{verdict} & \texttt{sufficient} or \texttt{insufficient} & Synthesizer & Stop-or-continue decision \\
    \hline\hline
  \end{tabular}
\end{table}

\section{System Overview}
\label{sec:methods-overview}

NanoMem has two stages. The first stage prepares the memory bank for retrieval.
It performs lightweight temporal normalization on session text, stores session
metadata, and indexes sessions or chunks in the memory database. This stage is
deliberately conservative: it makes obvious temporal expressions searchable but
does not try to precompute every future query interpretation.

The second stage runs when a query arrives. The query is normalized with respect
to query time, and a Planner produces retrieval cues. The Retriever searches the
memory database for candidate sessions. The Synthesizer reads the query,
retrieved sessions, temporal hints, and the current evidence state, then emits
reasoning, structured events, and a verdict. If the verdict is
\texttt{sufficient}, the loop stops and the final evidence is passed to the
answer model. If the verdict is \texttt{insufficient}, the Synthesizer's
reasoning and previously generated cues condition the next planning round.

This division gives each component a narrow responsibility. The Planner
generates search cues, the Retriever returns source material, the Synthesizer
constructs evidence and judges sufficiency, and the answer model performs final
answer generation. Chapter~\ref{chap:related-work} positioned this boundary as
the main distinction between NanoMem and answer-oriented reflective retrieval
systems.

\begin{figure}[H]
  \centering
  \fbox{\begin{minipage}{0.92\textwidth}
    \centering
    \vspace{0.8em}
    \textcolor{red}{TODO: Replace this placeholder with the planned NanoMem
    system architecture figure. It should show temporal normalization, memory
    database indexing, Planner, Retriever, Synthesizer, Temporal Evidence Pool,
    verdict loop, and downstream answerer, while marking the Synthesizer as the
    trainable component.}
    \vspace{0.8em}
  \end{minipage}}
  \caption{NanoMem system architecture. The memory module constructs structured
  evidence at query time and uses the verdict to decide whether to continue
  retrieval or stop.}
  \label{fig:methods-architecture}
\end{figure}

\section{Temporal Evidence Pool}
\label{sec:methods-evidence-pool}

The Temporal Evidence Pool $H_t$ is the state carried across retrieval rounds.
It is not a cache of raw retrieved sessions. It is a set of structured evidence
records:
\[
  h(e)=
  (\mathrm{text}(e),t_{\mathrm{event}}(e),t_{\mathrm{session}}(e),s(e)),
\]
where $\mathrm{text}(e)$ is a concise evidence statement,
$t_{\mathrm{event}}(e)$ is the event time, $t_{\mathrm{session}}(e)$ is the
session time, and $s(e)$ identifies the source session. The implementation
schema represents these fields with XML event text and the attributes
\texttt{session\_id}, \texttt{session\_time}, and \texttt{event\_time}.

The pool has two functions. Externally, it is the evidence passed to the answer
model. Internally, it is the working state of the search loop. A first retrieval
round may find an event that does not answer the query but reveals a time
window, entity, or causal predecessor. That information can then guide the next
round. In the running example, evidence about a promotion supplies the week in
which the system should search for reading-related evidence.

Each event is source grounded. A record cites one source session rather than a
set of sessions, which avoids hiding cross-session reasoning inside a single
claim. If multiple sessions are needed, the pool should contain multiple events,
each with its own source identifier. This design makes it easier to diagnose
whether an error came from retrieval, temporal grounding, synthesis, or final
answering.

\begin{figure}[H]
  \centering
  \fbox{\begin{minipage}{0.92\textwidth}
    \centering
    \vspace{0.8em}
    \textcolor{red}{TODO: Replace this placeholder with the planned Temporal
    Evidence Pool schema figure. It should show a query, two source sessions,
    and event records with \texttt{session\_id}, \texttt{session\_time},
    \texttt{event\_time}, event text, and verdict.}
    \vspace{0.8em}
  \end{minipage}}
  \caption{Record schema of the Temporal Evidence Pool. Separating event time
  from session time allows a mention in one session to become a temporal cue for
  searching another session.}
  \label{fig:methods-pool-schema}
\end{figure}

\section{Temporal Normalization and Hints}
\label{sec:methods-temporal-normalization}

NanoMem uses deterministic temporal normalization where possible. The purpose
is not to solve all temporal reasoning with rules, but to remove avoidable date
arithmetic from the learned Synthesizer. Relative expressions such as
\texttt{yesterday}, \texttt{last week}, \texttt{next month}, weekday names, and
several ``N days ago'' style phrases are resolved against a reference
timestamp. For sessions, the reference is the session timestamp. For the query,
the reference is the query timestamp.

The normalized text preserves the original wording while adding an absolute
span when a phrase can be resolved. For example, a phrase such as ``last week''
can be rewritten with the corresponding date interval. This makes temporal
information visible to lexical and embedding retrieval. If a session and a
query use different relative expressions for the same time interval, both can
still become searchable through a shared absolute span.

Temporal information has two distinct roles in the implementation. It remains
a hint for the Synthesizer: temporal spans are rendered in the prompt so that
the model can decide whether retrieved sessions support the query and whether a
new temporal cue should be searched in a later round. In the GRPO and
\texttt{search\_expr} retrieval path, however, parseable cue-level temporal
spans are also used as a targeted hard mention filter. The index search first
scores sessions with lexical and dense signals, with the default combined score
using normalized BM25 and dense embedding similarity. For the training rollout
path, temporal boost is disabled, reranking is disabled, and only the top seven
pre-reranker sessions per cue are kept by default. When a cue contains
temporal spans, \texttt{temporal\_filter\_mode=hard\_mention} then retains
candidate sessions whose session timestamp overlaps the cue span.

This choice is narrower than a universal session-time filter. The system does
not discard sessions by the original query time before cue generation, and the
Synthesizer still reasons over event time separately from session time. The
hard mention filter is used after the Planner has produced a cue with an
explicit temporal span, mainly to reduce candidate noise in temporal-bridge
rollouts. Its risk is recall loss: an answer-bearing event may be mentioned
outside the target interval even when its event time lies inside it. Chapter 4
therefore has to interpret retrieval failures with this trade-off in mind.

\begin{figure}[H]
  \centering
  \fbox{\begin{minipage}{0.92\textwidth}
    \centering
    \vspace{0.8em}
    \textcolor{red}{TODO: Replace this placeholder with the planned temporal
    normalization figure. It should show session text normalized against
    session time, query text normalized against query time, and Synthesizer
    event times becoming hints for later retrieval rounds.}
    \vspace{0.8em}
  \end{minipage}}
  \caption{Temporal normalization and hint generation. Deterministic parsing
  handles simple relative expressions, while event-time inference remains part
  of the Synthesizer's evidence construction task.}
  \label{fig:methods-temporal-normalization}
\end{figure}

\begin{table}[H]
  \centering
  \caption{Temporal-processing design choices.}
  \label{tab:methods-temporal-choices}
  \scriptsize
  \begin{tabular}{p{0.23\textwidth}p{0.25\textwidth}p{0.26\textwidth}p{0.16\textwidth}}
    \hline\hline
    Design issue & Adopted direction & Rationale & Main risk \\
    \hline
    Resolve simple relative expressions & Deterministic parser anchored to session or query timestamp & Reduces brittle date arithmetic for the Synthesizer & Rule coverage is incomplete \\
    Preserve original phrase and span & Keep phrase-level temporal metadata & Lets the model see both wording and resolved date & Extra prompt complexity \\
    Treat time as retrieval signal & Use temporal hints in the Synthesizer and cue-level \texttt{hard\_mention} filtering in GRPO retrieval & Reduces candidate noise after the Planner exposes a time span & May drop sessions whose mention time differs from event time \\
    Let event time enter later search & Use Synthesizer event times as loop state & Makes indirect evidence addressable & Wrong event time can drift search \\
    Unknown time & Preserve text and use \texttt{unknown} when needed & Avoids fabricating precision & Downstream reasoning receives weaker evidence \\
    \hline\hline
  \end{tabular}
\end{table}

\section{Planner--Retriever--Synthesizer Loop}
\label{sec:methods-loop}

At retrieval round $t$, the Planner produces a set of cues $C_t$. In the first
round, the cues are generated from the query. In later rounds, they are
generated from the query, previous cues, and the previous Synthesizer reasoning:
\[
  C_t =
  \begin{cases}
    P(q), & t=0,\\
    P(q,C_{<t},r_{t-1}), & t>0.
  \end{cases}
\]
The Planner is not responsible for answering. Its output is a retrieval plan in
the form of cue text and temporal phrase annotations.

The Retriever searches the indexed memory database for each cue. The NanoMem
codebase supports dense embeddings, query preprocessing, stateful retrieval,
index search, reranking, and budget-based pruning. In the GRPO rollout path,
the active L0 index search combines normalized BM25 and dense retrieval scores
with weights $0.25$ and $0.75$, respectively. The training path disables
temporal boost and reranking, applies cue-level hard mention filtering when
temporal spans are available, keeps the top seven pre-reranker sessions per
cue, merges results across cues, removes duplicate sessions, and trims complete
sessions until the Synthesizer input fits the token budget. These details are
method choices for the reported training/evaluation interface, not merely
appendix hyperparameters, because they define which evidence the Synthesizer can
see.

The Synthesizer then receives the query, temporal hints, retrieved sessions, and
the current state. It emits XML in the following shape:
\[
\begin{array}{l}
\texttt{<reasoning>}~r_t~\texttt{</reasoning>}\\
\texttt{<events>}~\mathcal{E}_t~\texttt{</events>}\\
\texttt{<verdict>}~v_t~\texttt{</verdict>.}
\end{array}
\]
Each event has a single \texttt{session\_id}, a \texttt{session\_time}, an
\texttt{event\_time}, and text. The verdict is binary:
\texttt{sufficient} stops the loop and passes the evidence to the answer model;
\texttt{insufficient} continues retrieval. This binary schema is the
method-level schema used by the current parser. Earlier design variants exposed
an \texttt{inferred} event attribute and an \texttt{indicative} verdict. In the
current method, the inference boundary is represented differently: event
records are treated as source-grounded claims tied to one cited session, while
cross-session or uncertain reasoning should remain in the \texttt{reasoning}
field or be decomposed into separate source-grounded events. Thus RQ1's concern
with inference markers is addressed as an attribution and reasoning-boundary
requirement rather than as a separate XML attribute in the final parser.

\[
  (r_t,\mathcal{E}_t,v_t)
  \sim \pi_\theta(\cdot \mid q,\widetilde{\mathcal{S}}_{\le t},H_t).
\]
If $v_t=\texttt{sufficient}$, the loop terminates. Otherwise, the reasoning
$r_t$ and prior cues are fed back into the Planner. This is what allows
evidence discovered in one round to become the search condition for a later
round.

\begin{figure}[H]
  \centering
  \fbox{\begin{minipage}{0.92\textwidth}
    \centering
    \vspace{0.8em}
    \textcolor{red}{TODO: Replace this placeholder with the planned dynamic
    loop figure. It should show two rounds: first finding a temporal anchor,
    then using that event time to retrieve the answer-bearing session.}
    \vspace{0.8em}
  \end{minipage}}
  \caption{Dynamic Planner--Retriever--Synthesizer loop. The loop is useful
  when the first evidence found is not the final answer but a bridge that makes
  the next retrieval query possible.}
  \label{fig:methods-loop}
\end{figure}

\begin{table}[H]
  \centering
  \caption{Synthesizer XML fields and constraints.}
  \label{tab:methods-xml}
  \scriptsize
  \begin{tabular}{p{0.18\textwidth}p{0.26\textwidth}p{0.25\textwidth}p{0.21\textwidth}}
    \hline\hline
    XML part & Required content & Meaning & Failure risk \\
    \hline
    \texttt{reasoning} & Non-empty text & Explains current evidence and missing information & Unhelpful reasoning gives poor cues for next round \\
    \texttt{events} & Zero or more event records & Structured evidence extracted from retrieved sessions & Missing events reduce answerability \\
    \texttt{event} attributes & \texttt{session\_id}, \texttt{session\_time}, \texttt{event\_time} & Source and temporal grounding & Missing or multi-session IDs break attribution \\
    Event text & Non-empty concise statement & Evidence statement for the answer model & Hallucinated or bloated evidence harms downstream answering \\
    \texttt{verdict} & \texttt{sufficient} or \texttt{insufficient} & Stop-or-continue control signal & Premature stop or unnecessary search \\
    \hline\hline
  \end{tabular}
\end{table}

\section{Training the Synthesizer with GRPO}
\label{sec:methods-training}

NanoMem trains the Synthesizer because it performs the hardest query-time
operation: extracting event-level evidence, assigning event times, and judging
whether the current evidence is enough. The Planner, Retriever, downstream
answer model, judge, and reference policy are treated as frozen environment
components during GRPO training. This keeps the credit assignment focused on
the Synthesizer. If both Planner and Synthesizer were updated at once, a failed
answer could be caused by a bad cue, a missed retrieval, a bad event, or a bad
verdict, making the training signal less interpretable.

For a query $q$, the training environment samples a group of rollout
trajectories. A trajectory contains one or more search rounds, online retrieval,
Synthesizer XML generation, verdict-based termination, and terminal answer
evaluation. Let
\[
  h_t = (q,\widetilde{\mathcal{S}}_{\le t},C_{\le t},H_t)
\]
be the retrieval context at round $t$. The Synthesizer action is
\[
  z_t=(r_t,\mathcal{E}_t,v_t)\sim\pi_\theta(\cdot\mid h_t).
\]
Only Synthesizer-generated XML tokens are part of the policy action. Planner
outputs, retrieved sessions, prompts, and other environment text are masked out
from the policy loss. The implementation records this distinction with a
response mask over rollout tokens.

NanoMem uses Group Relative Policy Optimization (GRPO) \cite{shao2024deepseekmath}.
For each query, a group of $G$ trajectories
$\mathcal{T}=\{\tau^i\}_{i=1}^G$ receives scalar rewards. The advantage for a
trajectory is normalized within the group:
\[
  \hat{A}^i =
  \frac{R(\tau^i)-\mathrm{mean}_{j=1}^{G}R(\tau^j)}
  {\mathrm{std}_{j=1}^{G}R(\tau^j)+\epsilon}.
\]
The objective applies the clipped policy update to the unmasked Synthesizer
tokens and includes a KL penalty against the reference policy. This makes the
training closer to multi-step agent learning than to single-turn instruction
tuning: the model is rewarded for producing evidence that works inside a search
environment.

\begin{figure}[H]
  \centering
  \fbox{\begin{minipage}{0.92\textwidth}
    \centering
    \vspace{0.8em}
    \textcolor{red}{TODO: Replace this placeholder with the planned GRPO rollout
    and token-masking figure. It should show several trajectories for the same
    query, environment tokens masked out, and Synthesizer XML tokens optimized.}
    \vspace{0.8em}
  \end{minipage}}
  \caption{GRPO rollout and token masking. NanoMem optimizes the Synthesizer's
  XML output while treating planning, retrieval, and answer evaluation as the
  environment.}
  \label{fig:methods-grpo}
\end{figure}

\section{Reward Design}
\label{sec:methods-reward}

The reward is designed for evidence construction rather than direct answer
generation. The checked GRPO reward implementation combines four weighted
signals: format validity, event faithfulness, outcome usefulness, and length.
The verdict remains central to the method, but it is not a separate additive
reward component in the active reward total. Instead, it controls whether a
trajectory is terminal and therefore whether terminal outcome evaluation is
available.

\textbf{Format reward.} The Synthesizer output must be parseable XML with
\texttt{reasoning}, \texttt{events}, required event attributes, non-empty event
text, and a valid binary verdict. Format validity is necessary because the rest
of the pipeline depends on structured evidence. A malformed output cannot be
reliably attributed, evaluated, or passed downstream.

\textbf{Event faithfulness reward.} Each emitted event is supposed to be
directly supported by its cited session. The faithfulness component checks
event text against the corresponding session with a frozen verifier and rewards
the proportion of supported events. This signal is needed because a compact
evidence pool is useful only if compression preserves attribution. Since the
current parser does not expose an \texttt{inferred} field, the method-level
contract is stricter than the earlier design: every emitted event should be
source-grounded. Evidence that depends on cross-session inference should be
expressed through separate cited events and the \texttt{reasoning} field rather
than as a single unsupported event.

\textbf{Outcome reward.} When the loop terminates, the evidence is provided to
a frozen answer model. A judge compares the answer to the gold answer. This
reward connects the memory artifact to task usefulness, but it does not make
NanoMem a final answer agent. The trained policy is still the evidence
Synthesizer.

\textbf{Length reward.} Compactness is part of the evidence-provider objective:
the answer model should receive enough evidence, not an unbounded transcript.
The length reward encourages shorter Synthesizer outputs, usually conditional
on the output being useful. It must not dominate outcome or verdict quality,
because a trivially short but uninformative evidence pool is not useful memory.

A representative total reward is
\[
  R(\tau)=
  w_f R_{\mathrm{fmt}}(\tau)
  + w_s R_{\mathrm{faith}}(\tau)
  + w_o R_{\mathrm{out}}(\tau)
  + w_l R_{\mathrm{len}}(\tau),
\]
with default MVP weights $w_f=0.2$, $w_s=0.2$, $w_o=0.5$, and $w_l=0.1$ in the
configuration artifact. The intended priority is outcome usefulness first, then
parseability and source support, then compactness. Verdict quality is evaluated
through loop behavior and downstream outcome: a premature \texttt{sufficient}
verdict exposes incomplete evidence to the answer model, while repeated
\texttt{insufficient} verdicts can terminate by maximum-round or no-new-evidence
conditions and receive no terminal outcome reward.

\begin{figure}[H]
  \centering
  \fbox{\begin{minipage}{0.92\textwidth}
    \centering
    \vspace{0.8em}
    \textcolor{red}{TODO: Replace this placeholder with the planned reward
    signal figure. It should mark format and faithfulness on the Synthesizer
    artifact, outcome at terminal answer evaluation, compactness on Synthesizer
    output length, and the verdict as the terminal gate.}
    \vspace{0.8em}
  \end{minipage}}
  \caption{Reward signals in a NanoMem trajectory. Format and faithfulness
  supervise the structured evidence artifact, while outcome and length connect
  terminal evidence to downstream usefulness and compactness. The verdict gates
  whether terminal outcome evaluation is reached.}
  \label{fig:methods-rewards}
\end{figure}

\begin{table}[H]
  \centering
  \caption{Reward components and their roles.}
  \label{tab:methods-rewards}
  \scriptsize
  \begin{tabular}{p{0.18\textwidth}p{0.21\textwidth}p{0.19\textwidth}p{0.23\textwidth}p{0.09\textwidth}}
    \hline\hline
    Component & Measures & Timing & Failure mode addressed & Risk \\
    \hline
    Format & XML parseability and required fields & Each Synthesizer output & Unusable evidence artifact & Models may learn shallow formatting \\
    Faithfulness & Direct support between event text and cited session & Parsed events in the terminal reward path & Unsupported compressed evidence & Requires parser and verifier assumptions to stay aligned \\
    Outcome & Answer correctness from terminal evidence & Terminal only & Evidence that is well formed but not useful & Can hide which component failed \\
    Length & Compactness of Synthesizer output & Usually terminal or trajectory aggregate & Passing bloated evidence to answer model & Must not reward empty evidence \\
    \hline\hline
  \end{tabular}
\end{table}

\section{ChainMem Benchmark Construction}
\label{sec:methods-chainmem}

LoCoMo and LongMemEval evaluate long-term memory and multi-session reasoning
\cite{maharana2024locomo,wu2024longmemeval}, but many examples remain solvable
when the answer-bearing session is directly retrievable. NanoMem therefore uses
ChainMem as a diagnostic benchmark for indirectly addressable evidence. The
benchmark targets cases where the system must first retrieve or synthesize a
bridge before it can search for the final evidence.

Each ChainMem example contains a source event, a bridge relation, and a
destination event. The query makes the source event locatable, but the answer
is stored in the destination event. A valid system must recover the source
event time, derive the target window or condition, retrieve the destination
event, and then synthesize evidence for the answer model. Distractors are
included so that surface matching alone is not enough.

The construction process can be understood as event-graph-first generation.
First, a persona or theme and source event are sampled. Second, a bridge scope
and temporal offset determine the target window. Third, a destination event and
answer are placed inside or around that window. Fourth, hard distractors are
added, such as near-window events, answer aliases, or semantically similar
activities. Fifth, the event graph is realized as multi-turn sessions with
relative or indirect temporal expressions. Finally, automatic checks filter
invalid generations, temporal inconsistencies, split leakage, and non-unique
answers.

ChainMem is not intended to cover all long-term memory behavior. It is a stress
test for temporal-bridge retrieval. Public long-memory benchmarks remain
necessary for evaluating generality, while ChainMem isolates the failure mode
that motivates the iterative evidence-state design.

\begin{figure}[H]
  \centering
  \fbox{\begin{minipage}{0.92\textwidth}
    \centering
    \vspace{0.8em}
    \textcolor{red}{TODO: Replace this placeholder with the planned ChainMem
    construction figure. It should show source event, bridge window,
    destination event, distractors, and realization into sessions.}
    \vspace{0.8em}
  \end{minipage}}
  \caption{ChainMem construction. The benchmark creates questions whose final
  evidence becomes retrievable only after an intermediate temporal bridge has
  been recovered.}
  \label{fig:methods-chainmem}
\end{figure}

\begin{table}[H]
  \centering
  \caption{ChainMem construction variables and validation checks.}
  \label{tab:methods-chainmem}
  \scriptsize
  \begin{tabular}{p{0.24\textwidth}p{0.34\textwidth}p{0.32\textwidth}}
    \hline\hline
    Variable or check & Meaning & Why it matters \\
    \hline
    Source event & Event named or implied by the query & Provides the first retrieval target \\
    Bridge window & Temporal interval derived from the source event & Creates the hidden retrieval cue \\
    Destination event & Event containing the answer & Tests whether the system can continue search \\
    Distractor mode & Similar events, near times, aliases, or confusable activities & Prevents trivial surface matching \\
    Temporal style & Relative, calendar-based, or indirect phrasing & Tests event-time grounding \\
    Temporal consistency check & Ensures generated dates and intervals agree & Prevents invalid samples \\
    Answer uniqueness check & Ensures one answer is supported & Makes automatic evaluation meaningful \\
    Split leakage check & Prevents train/evaluation contamination & Protects benchmark validity \\
    \hline\hline
  \end{tabular}
\end{table}

\section{Implementation Boundary}
\label{sec:methods-implementation}

The NanoMem implementation is organized around a service-style memory pipeline.
Core runtime code lives under \texttt{src/nanomem}. The time utilities implement
relative temporal expression resolution and canonicalization. Storage and vector
store modules provide memory persistence and retrieval indexes. Operator modules
implement query preprocessing, read-path retrieval, compression, extraction,
and write-path behavior. The server and service modules expose the memory
pipeline for evaluation and training.

Training code is organized under \texttt{training/grpo}. The schema parser
defines the final Synthesizer event fields and valid verdicts. The rollout code
connects Planner prompts, online retrieval, Synthesizer generation, token
masking, and reward computation inside a GRPO agent loop. Evaluation code under
\texttt{evals} provides benchmark-specific runners for LoCoMo, LongMemEval,
ChainMem, and related baselines.

The implementation boundary used in this chapter is therefore concrete. The
reported GRPO retrieval path uses BM25/dense index search with the $0.25/0.75$
combined score, cue-level \texttt{hard\_mention} temporal filtering when cue
spans exist, disabled reranking, top-seven pre-reranker session retention per
cue, complete-session token-budget pruning, and up to three retrieval rounds by
default. The active reward total uses format, faithfulness, length, and outcome
weights. Because the current event parser no longer exposes the older
\texttt{inferred} attribute, any reported training run must keep the
faithfulness reward path synchronized with the parser and treat emitted events
as directly source-grounded claims. Chapter 4 should still report run-specific
model identities, data splits, group size, checkpoint selection, and any
configuration overrides, but the retrieval and reward behavior above is part of
the method definition.

\begin{table}[H]
  \centering
  \caption{Mapping from method components to implementation areas.}
  \label{tab:methods-implementation}
  \scriptsize
  \begin{tabular}{p{0.22\textwidth}p{0.27\textwidth}p{0.20\textwidth}p{0.21\textwidth}}
    \hline\hline
    Method component & Main implementation area & Input/output & Method-relevant setting or check \\
    \hline
    Temporal normalization & \texttt{src/nanomem/time\_utils.py} & Text plus timestamp to resolved spans & Phrase coverage and date format \\
    Query cues & \texttt{src/nanomem/operators/query.py}; prompt cue schema & Query to structured cues & Planner cues are enriched with temporal spans before retrieval \\
    Retrieval & \texttt{src/nanomem/operators/read\_path.py}; vector stores & Cues to candidate sessions & BM25/dense $0.25/0.75$, \texttt{hard\_mention}, reranker off, top-seven per cue \\
    Synthesizer schema & \texttt{training/grpo/trainer/schemas.py} & XML to structured events and verdict & \texttt{session\_id}, \texttt{session\_time}, \texttt{event\_time}, text, binary verdict \\
    GRPO rollout & \texttt{training/grpo/trainer/rollout.py} & Query trajectory to masked policy tokens & Maximum three rounds by default; terminal on \texttt{sufficient} \\
    Reward pipeline & \texttt{training/grpo/trainer/rewards.py}; reward weights config & XML and terminal answer to scalar reward & Format, faithfulness, length, and outcome components \\
    Benchmarks & \texttt{evals/locomo}, \texttt{evals/longmemeval}, \texttt{evals/chainmem} & Dataset records to evaluation results & Split sizes and filtering rules \\
    \hline\hline
  \end{tabular}
\end{table}

\section{Chapter Summary}
\label{sec:methods-summary}

This chapter defined NanoMem as an evidence-state construction system. The
Temporal Evidence Pool answers RQ1 by representing long-term memory as
source-grounded events with separate session-time and event-time fields. The
Planner--Retriever--Synthesizer loop answers RQ2 by allowing evidence found in
one round to guide retrieval in the next round. GRPO training and the reward
design align the Synthesizer with parseable, compact, useful, and
sufficiency-aware evidence construction. Chapter 4 evaluates
whether these design choices improve answer quality, temporal grounding,
compactness, and hidden-dependency retrieval on public memory benchmarks and
ChainMem.
